\section{Spiritual Man and the Tyrant}

\begin{quotex}
It is one's psychological type which from the outset determines and limits a person's judgment
\flright{\textsc{Carl Jung}}
\end{quotex}


In a book review of \emph{Monas Hieroglyphica} (or \emph{Hieroglyphic Monad}) by \textbf{John Dee}, \textbf{Valentin Tomberg} describes the two life paths described by Dee in the preface, incorporating into it the seven-year cycles of life described by \textbf{Rudolf Steiner}. The first stages of life are the same for both paths, as interpreted by Tomberg. I have added the sacraments associated with the beginning of each stage, even if their esoteric significance has apparently long been forgotten. Since Tomberg does not mention Steiner's description of the stages, I also included them.

\begin{table}[h]\small\centering
\begin{tabular}{*{4}{p{.2\textwidth}}}\toprule

Stage & Age & Sacrament & Description \\\midrule

Infancy \& Childhood & Up to 7 & Baptism & The psychic forces transform the body of the child from one that was inherited from the parents, to one that represents the full personality of the child. \\

Youth & 7 to 14 & Communion/Age of Reason & The child's imagination and feeling life takes centre stage \\

Adolescence & 14 to 21 & Confirmation & The higher mind of the adolescent takes root \\

Young adulthood & 21 to 28 & ~ & Sentient soul \\\bottomrule
\end{tabular}
\end{table}

During young adulthood, and even adolescence, a choice can be made:

\begin{itemize}


\item The path of earth or worldliness

\item The path of water or wisdom
\end{itemize}


The first path leads to the \textbf{Tyrant} ending in the Abyss; the second leads to the \textbf{Spiritual Man} (\emph{pneumatikos}), ending in Fire. These are the stages on each path, including Tomberg's addition of the ages.

\begin{table}[h]\small\centering
\begin{tabular}{ccc}
\toprule
Tyrant & Age & Spiritual Man \\\midrule

Earth & 35 & Water \\

Trouble & 42 & Philosopher \\

Deception & 49 & Sage \\

Power & 56 & Adeptship\\\bottomrule
\end{tabular}
\end{table}

Tomberg interprets this to indicate that Dee has a firm understanding of the stages of life. However, in point of fact, few people make that choice of all, becoming neither tyrant or philosopher; or else they only go partway along the path. Most people, on the other hand, seem to wander aimlessly through life, oblivious to any higher calling.

\paragraph{IQ and Intelligence}


\begin{quotex}
If you're so smart, why aren't you rich?
\flright{\textsc{old saying}}
\end{quotex}


\textbf{Nassim Taleb} recently tweeted something about IQ tests:


\begin{quotex}
IQ measures an inferior form of intelligence, stripped of 2nd order effects, meant to select paper shufflers, obedient IYIs intelligent yet idiots.
\end{quotex}


Taleb is normally reliable, but his desire to be the gadfly and curmudgeon --- i.e., to be outrageous --- might get the better of him. Apparently he believes that ``higher'' intelligence applies solely to being successful in the world. Actually, success in the world does not required an exceptionally high IQ, or else the human race would have failed millennia ago.

On the contrary, the highly intelligent are less interested in success by worldly standards and are more attracted to abstract thinking. In other words, they take the water path rather than the earth path. Thus, even if they are not highly financially successful, they are hardly obedient paper shufflers. Instead of IQ, perhaps a more objective test can be used. Here are some proposals, for example, based on academic courses:

\begin{itemize}


\item Pass a course in logic. I was in a class that started with 25 students and ended with just 3 or us. The philosophy, political science, pre-law students had all dropped out.

\item Read and understand one of Plato's dialogs.

\item Pass an economics class. This class forced many business majors to switch majors.

\item Pass organic chemistry. This is often the bane of pre-med students.

\item Pass a physics class. Few students even attempt it.

\item Pass a class in abstract mathematics. I used to teach algebra to elementary education majors, many of whom were taking the course for the second or third time. That was torture.

\item Read and comment on a few sonnets.

\item Read and comment on selected works of fiction.
\end{itemize}


Obviously, you can do well in the world without taking any of those classes. That does not mean that intelligence does not exist.

\paragraph{Twofold Reason}


\textbf{Aristotle} identified two types of reason: Creative and Passive.

These are described as:

\begin{itemize}


\item \textbf{Creative Mind}: Always in full possession of all forms or ideas.

\item \textbf{Passive Mind}: Only has potential possession of them and which can only rise to actual possession through the inspiration of the creative mind.
\end{itemize}


Most students who are considered smart in school work through the passive mind. That is, they are really competent at absorbing and repeating the ideas taught by their teachers. Of course, in most cases the teachers themselves only have a passive mind. Such students are convinced they are intelligent because they can absorb the propaganda repeated by the powerful.

However, the ideas themselves come from those with creative minds. They can resist propaganda due to their better critical thinking skills. They may be considered to be oddballs and may fail in situations which required conformity.

It is true that IQ may not capture this distinction.

\paragraph{Things and Concepts}


\begin{quotex}
I \emph{must}, before I die, find some way to say the essential thing that is in me, that I have never said yet --- a thing that is not love or hate or pity or scorn, but the very breath of life, of development, fierce and coming from far away, bringing into human life the vastness and the fearful passionless force of non-human things.
\flright{\textsc{Bertrand Russell}, \emph{My Philosophical Development}}
\end{quotex}


Bertrand Russell separates terms into two types: things and concepts. Things refer to individuals and concepts to classes. For Russell, these are merely logical categories, so there is no judgment about the reality of concepts.

In reality, the thing, or percept, refers to existence---or that the thing exists. Concepts refer to essences --- what a thing is. The dominant position today is that concepts are not real. Rather, they are just nominal concepts --- or social constructs in today's language --- and do not refer to anything transcendental.

The ability to ``see'' the concept in the percept is an ability not available to everyone (although it can be learned). This is a form of theophany since the concepts are really ideas in the Mind of God. Borrowing form \textbf{Carl Jung}, with some poetic license, we can distinguish two different perceiving types: the intuitive and the sensitive.

The sensitive type is dominated by sensation, i.e., he experiences the thing through the senses. The intuitive type, on the other hand, can ``see'' the concept in the thing, which is a direct knowing of the idea. The intuitive type can certainly experience the thing via the senses, just as the sensitive type does, but he also experiences reality more deeply. Specifically, he also sees the idea instantiated by the thing. \textbf{John Findlay} affirms this:


\begin{quotex}
It seems plain that our human world is a world in which innumerable ideas, meanings, facts, principles, constructions, hypotheses, laws, images and ideals are as essential a part of the landscape as are the concrete bodies and thinking persons around which they cluster, and above which they float. They constitute a universal world of rational mind in which all thinking persons share, whatever the limitations of their immediate, sensuous viewpoint.
\end{quotex}


\paragraph{Instantiation}


Statistics were unknown to the ancients and are very poorly understood today, even by the educated. Qualities follow a statistical distribution when they become instantiated in a population. This divergence leads many to think that the quality is not real or is irrelevant. Findlay, in his reworking of Plato, makes this clear:

\begin{quotex}
Probability is an irremovable feature of the human cave, and as such it resists reduction to anything factual, anything ascertainable and discoverable, anything that is could be merely there.
\end{quotex}


That means that the ideal of necessary connections is not always met. When a concept or quality manifests itself in existence, it will follow a statistical distribution of some sort (otherwise, there would be a hidden law or principle). So for intelligence, there will be a range from the highly to the average and to low intelligence. This does not mean that intelligence does not exist. It also precludes the possibility of any notion of egalitarianism in intelligence.

The instantiation of a divine idea requires sacrifice and limitation. Otherwise, the essence and existence of finite beings would be identical. However, the identity of essence and existence is possible only for God.

\textit{References}

\textit{The Hieroglyphic Monad}\footnote{\url{http://www.esotericarchives.com/dee/monad.htm}}

\textsc{Valentin Tomberg}, \emph{Russian Spirituality}

\textit{Stages of Life according to Rudolf Steiner}\footnote{\url{http://www.institute4learning.com/2012/08/07/the-stages-of-life-according-to-rudolf-steiner/}}

\textsc{John Findlay}, \emph{The Discipline of the Cave}

\textsc{Bertrand Russell}, \emph{Principles of Mathematics}

\textsc{Carl Jung}, \emph{Psychological Types}



\flrightit{Posted on 2018-12-27 by Cologero}

\begin{center}* * *\end{center}

\begin{footnotesize}
\begin{sffamily}
\texttt{James on 2018-12-28 at 08:42 said:}

It is interesting to think about the ``turning point'' (midway point) of the human journey at thirty five leading either to earth or water . I had written tangentially about this topic in my Dante book (still unpublished) because I had taken note in one of my chapters about how Dante (at the turning point age) finds himself on the same WATER journey as his predecessor poets .

Homer spoke of Odysseus crossing the sea from Troy back to his home and Virgil speaks of another departure from Troy to the eventual founding of Rome via another sea voyage . What Dante accomplishes , if I might be so bold to propose , is the synthesis of the earth and water voyages . He demonstrates that the esoteric waters the Christian (of the post-ancient age) must cross is the elemental dissolution of sin itself . He traverses the ocean of sin in order to find his true home by passing the sphere of fire above Purgatory as he ascends to Heaven . His accomplishment is through the assimilation-without-contamination or the understanding-without-touching of evil itself--of the path of Earth . In fact , quite ``literally'' , as he passes through the Earth and frozen water , he finds his way out of such a place even after meeting the ultimate Tyrant and ultimate Abyss at the bottom .

This is not so dissimilar to the various comments mentioned among friends about how resisting the devil can unlock a willing servant . To pass through the path of Earth with the understanding of Water --- in a sense as if to baptize ourselves (we being of dust and clay -- Earth) --- one can achieve the ``other path'' that is not the path of the exoteric (``the straight path had been lost'') or the path of martyrdom (pressing on despite the presence of the three beasts) . To use Satan himself as a willing ladder to reach Purgatory on the other side is part of the lesson I am trying to meditate upon . I wish to spend some time confirming in meditation if Dante truly does attempt this synthesis of Earth and Water both by demonstrating its various points of union (such as in the negative synthesis of the frozen water in Cocytus) . Perhaps it is the human person who is the representative of Earth that submits to baptism that gives a clue . Still , it is only a nascent idea in my mind that I probably should have kept to myself in order to confirm first with thought and experience , but I couldn't help but find the connections compelling . There must have been a reason that Steiner , Tomberg , and others could find meaning at the ``midway'' point of life's journey and how the elements associated pair up perfectly with what Dante was attempting to sing about .

\hfill

\texttt{Han Fei on 2018-12-29 at 12:05 said:}

I have a question. If a person knows himself, or at least beginning to do so, what kind of opinion should he have of himself? What is the extent of the ability that man is capable of? The reason why this may be relevant is because Taleb seems to consider intelligence as having to do with success, that is to say to the extent of some activity we perform in our life.

If I say am intelligent, then intelligent for what? If I say that I am free (as in self-agent) for what sake does this freedom hold? The same question applies to women as well.

\hfill

\texttt{Cologero on 2018-12-31 at 10:12 said:}

\textit{Union of Earth and Water by Rubens}\footnote{\url{https://www.freeart.com/gallery/r/rubens/rubens95.html}}

\hfill

\texttt{Lyon on 2019-01-08 at 21:22 said:}

Brilliant. Yet again. Time to do a deeper dive into Valentin Tomberg and Rudolf Steiner. This may take a few weeks or decades.

\hfill

\end{sffamily}
\end{footnotesize}

